%% LabelFusion-TS on FOMC stance — internal draft for co-author discussion.
%% Build: latexmk -pdf acl_latex.tex
\documentclass[11pt]{article}
\usepackage[review]{acl}
%% ===========================================================================
%%  preamble.tex — shared design system
%%
%%  Loaded by BOTH acl_latex.tex (the paper) and design_kit.tex (the pattern
%%  gallery), so anything defined here is guaranteed to look the same in both.
%%  Never put content here; content lives in sections/.
%%
%%  Hard constraints this file respects (see formatting.md, paper/guideline.md):
%%   * acl.sty is official and UNMODIFIED — all of our styling happens here.
%%   * acl.sty already loads: geometry, caption (font=10pt), natbib, hyperref,
%%     and (in review mode) lineno. Loading those again with options would
%%     raise an option clash, so we only ever \captionsetup{...} them.
%%   * Body font must stay Times at 11pt with the style's leading; we adjust
%%     float/caption spacing only, never \baselineskip or the type area.
%% ===========================================================================


%% ---------------------------------------------------------------------------
%%  0. Standard includes (verbatim from the official acl_latex.tex)
%% ---------------------------------------------------------------------------

\usepackage{times}
\usepackage{latexsym}
\usepackage[T1]{fontenc}     % correct rendering/hyphenation of accented words
\usepackage[utf8]{inputenc}
\usepackage{microtype}       % better line breaking; saves a surprising amount
\usepackage{inconsolata}     % typewriter face for code and prompts


%% ---------------------------------------------------------------------------
%%  1. Packages
%% ---------------------------------------------------------------------------

% -- tables ----------------------------------------------------------------
\usepackage{booktabs}       % \toprule \midrule \bottomrule — the only rules we use
\usepackage{multirow}       % vertically merged cells
\usepackage{makecell}       % line breaks inside a cell: \makecell{a\\b}
\usepackage{tabularx}       % X columns that soak up the remaining width
\usepackage{threeparttable} % table notes that inherit the table's width
\usepackage{supertabular}   % the ONLY long-table package that breaks across
                            % columns in two-column mode (longtable cannot)
\usepackage{siunitx}        % S columns for decimal-aligned numbers.
                            % Only v2+v3-portable syntax is used: S[table-format=...]
                            % and \num{...}. Do not add \sisetup keys that were
                            % renamed in siunitx v3 (e.g. table-number-alignment).
\usepackage{adjustbox}      % max width=\linewidth — shrink-to-fit, see \fitbox

% -- figures ---------------------------------------------------------------
\usepackage{graphicx}
\usepackage{subcaption}     % subfigures; must come after caption (acl.sty loaded it)

% -- algorithms ------------------------------------------------------------
\usepackage{algorithm}      % the float
\usepackage{algpseudocode}  % the pseudocode body

% -- code and boxes --------------------------------------------------------
\usepackage{listings}
\usepackage{xcolor}
\usepackage{tcolorbox}
\tcbuselibrary{breakable,skins,listings}

% -- symbols, math, lists --------------------------------------------------
\usepackage{amsmath}
\usepackage{amssymb}
\usepackage{pifont}         % \ding{51} checkmark for the vote profiles
\usepackage{enumitem}

% -- cross references (LAST: cleveref must see hyperref and all float types) --
\usepackage[capitalise,nameinlink]{cleveref}


%% ---------------------------------------------------------------------------
%%  2. Spacing
%%
%%  ACL papers are page-budget-bound (8 pages of content). These are the only
%%  legitimate savings: float separation and caption skips. Everything here is
%%  deliberately mild — if a draft overruns, cut text, do not shrink type.
%% ---------------------------------------------------------------------------

\setlength{\textfloatsep}{10pt plus 2pt minus 2pt}   % float <-> text, top/bottom
\setlength{\floatsep}{10pt plus 2pt minus 2pt}       % float <-> float
\setlength{\intextsep}{10pt plus 2pt minus 2pt}      % around [h] floats
\setlength{\dblfloatsep}{10pt plus 2pt minus 2pt}    % same, spanning floats
\setlength{\dbltextfloatsep}{12pt plus 2pt minus 2pt}

\captionsetup{skip=5pt}                              % caption <-> content
\captionsetup[sub]{font=small,skip=3pt}

% Table typography: booktabs wants no vertical rules and a little air.
\setlength{\tabcolsep}{4.5pt}
\renewcommand{\arraystretch}{1.06}

% Compact lists (ACL drafts overflow on list padding more than anywhere else).
\setlist[itemize]{leftmargin=1.2em,topsep=2pt,itemsep=1pt,parsep=0pt}
\setlist[enumerate]{leftmargin=1.4em,topsep=2pt,itemsep=1pt,parsep=0pt}
\setlist[description]{leftmargin=0pt,topsep=2pt,itemsep=1pt,parsep=0pt,style=nextline}

% Float placement discipline: ACL floats belong at the top (or bottom) of a
% column/page, never mid-text. Defaults nudge LaTeX in that direction; write
% [t] explicitly on every float anyway.
\renewcommand{\topfraction}{0.9}
\renewcommand{\bottomfraction}{0.7}
\renewcommand{\textfraction}{0.1}
\renewcommand{\floatpagefraction}{0.75}
\setcounter{topnumber}{3}
\setcounter{dbltopnumber}{3}


%% ---------------------------------------------------------------------------
%%  3. Colours
%%
%%  One restrained palette, reused by boxes, code, and figures. Every colour
%%  must survive greyscale printing and be distinguishable to colour-blind
%%  readers, so hue never carries meaning alone (boxes also differ by title).
%% ---------------------------------------------------------------------------

\definecolor{cpInk}{HTML}{1A1A1A}       % text on light fills
\definecolor{cpRule}{HTML}{9AA0A6}      % box frames, hairlines
\definecolor{cpWash}{HTML}{F4F5F7}      % neutral fill (system prompts, code)
\definecolor{cpBlue}{HTML}{2B5C8A}      % user/input, primary accent
\definecolor{cpGreen}{HTML}{2E6B4F}     % model output
\definecolor{cpAmber}{HTML}{8A6A1F}     % judge / rubric
\definecolor{cpPlum}{HTML}{6B3F7A}      % examples, retrieved evidence
\definecolor{cpRed}{HTML}{A33A32}       % failures, drafting notes


%% ---------------------------------------------------------------------------
%%  4. Tables — the three patterns
%%
%%  (a) column-width table   : table  + \fitbox   -> shrinks only if needed
%%  (b) page-width table     : table* + \fitbox   -> spans both columns
%%  (c) multi-page table     : \begin{cptablelong} (supertabular), appendix only
%%
%%  Rules for all three: caption ABOVE the tabular, booktabs rules only,
%%  no vertical lines, units in the header, \best to mark winners.
%% ---------------------------------------------------------------------------

% Shrink-to-fit wrapper. \linewidth resolves to the column inside `table` and
% to the full text width inside `table*`, so the SAME wrapper serves both and
% a table can never overflow its container regardless of content. It only ever
% scales down (max width), so small tables keep the body font size.
\newenvironment{fitbox}
  {\begin{adjustbox}{max width=\linewidth,center}}
  {\end{adjustbox}}

% Font size for table bodies. Change here, not per table.
\newcommand{\tabfont}{\small}

% Result marking (bold = best, underline = runner-up; both readable in grey).
\newcommand{\best}[1]{\textbf{#1}}
\newcommand{\second}[1]{\underline{#1}}
\newcommand{\na}{\textendash}                    % measured, not applicable
\newcommand{\nr}{$\emptyset$}                    % not run (matches vote profile)

% Significance markers for comparison tables.
\newcommand{\sigone}{\textsuperscript{*}}
\newcommand{\sigtwo}{\textsuperscript{**}}
\newcommand{\sigthree}{\textsuperscript{***}}

% Header cell that wraps in a narrow column: \hcell{false alarm\\rate}
\newcommand{\hcell}[1]{\makecell[bc]{#1}}

% Group heading spanning several columns, with a rule underneath.
\newcommand{\gcol}[2]{\multicolumn{#1}{c}{#2}}

% Multi-page table in a single column (appendix). supertabular, not longtable:
% longtable is illegal in two-column mode. Caption goes in \topcaption before.
\newenvironment{cptablelong}[1]
  {\tabfont\begin{supertabular}{#1}}
  {\end{supertabular}}


%% ---------------------------------------------------------------------------
%%  5. Figures
%%
%%  Export conventions (keep figures honest and legible at print size):
%%   * vector PDF, never PNG, unless the plot has raster content
%%   * single column: 3.15 in wide  ·  spanning: 6.5 in wide
%%   * label/tick font 8 pt so it matches this caption size after \includegraphics
%%     at width=\columnwidth (i.e. do NOT scale figures in LaTeX)
%%   * no in-figure titles — the caption is the title
%% ---------------------------------------------------------------------------

% Grey placeholder that reserves the real figure's footprint. Lets us lay out
% and page-count before any plot exists: \figph{\columnwidth}{4.2cm}{caption id}
\newcommand{\figph}[3]{%
  \setlength{\fboxsep}{0pt}%
  \fcolorbox{cpRule}{cpWash}{%
    \parbox[c][#2][c]{#1}{\centering\small\color{cpRule}\ttfamily #3}}}


%% ---------------------------------------------------------------------------
%%  6. Algorithms
%%
%%  (a) single column : \begin{algorithm}[t]
%%  (b) page width    : \begin{algorithm*}[t]  (float, cannot break)
%%  (c) unbounded     : \begin{cpalgorithm}    (not a float: breaks across
%%                       columns and pages, for long appendix pseudocode)
%% ---------------------------------------------------------------------------

\floatname{algorithm}{Algorithm}
\algrenewcommand\algorithmicrequire{\textbf{Input:}}
\algrenewcommand\algorithmicensure{\textbf{Output:}}
\algrenewcommand{\algorithmiccomment}[1]{\hfill\small\textcolor{cpRule}{\# #1}}
\algnewcommand{\LineComment}[1]{\State\small\textcolor{cpRule}{\# #1}}
\newcommand{\algfont}{\small}

% Non-floating algorithm: identical look, but the body may split. Numbering
% shares the `algorithm` counter, so references stay consistent with (a)/(b).
\newenvironment{cpalgorithm}[1][]
  {\par\addvspace{8pt}\refstepcounter{algorithm}%
   \noindent\rule{\linewidth}{0.8pt}\par\vspace{2pt}%
   \noindent\textbf{\algorithmname~\thealgorithm}%
   \ifx\relax#1\relax\else\quad #1\fi\par\vspace{1pt}%
   \noindent\rule{\linewidth}{0.4pt}\par\vspace{2pt}\algfont}
  {\par\vspace{2pt}\noindent\rule{\linewidth}{0.8pt}\par\addvspace{8pt}}
% \algorithmname is defined by \floatname above; alias for the environment:
\providecommand{\algorithmname}{Algorithm}


%% ---------------------------------------------------------------------------
%%  7. Code
%%
%%  Two shapes only:
%%   \begin{pycode}[caption]  — verbatim Python, breaks across columns/pages
%%   \begin{pycodenum}[...]   — same with line numbers (for text that refers
%%                              to specific lines)
%%  Page-width code must sit inside a figure* float (floats cannot break).
%%  Inline: \code{EnsembleDetector} for identifiers in running text.
%% ---------------------------------------------------------------------------

\newcommand{\code}[1]{\texttt{\small #1}}

\lstdefinestyle{cppython}{
  language=Python,
  basicstyle=\scriptsize\ttfamily\color{cpInk},
  keywordstyle=\color{cpBlue}\bfseries,
  commentstyle=\color{cpGreen},
  stringstyle=\color{cpPlum},
  identifierstyle=\color{cpInk},
  showstringspaces=false,
  breaklines=true,
  breakatwhitespace=false,
  postbreak=\mbox{\textcolor{cpRule}{$\hookrightarrow$}\space},
  tabsize=4,
  columns=fullflexible,
  keepspaces=true,
  upquote=true,
  literate={-}{-}1 {'}{'}1,       % keep hyphens/quotes copy-pasteable
}
\lstdefinestyle{cpjson}{
  basicstyle=\scriptsize\ttfamily\color{cpInk},
  stringstyle=\color{cpPlum},
  showstringspaces=false,
  breaklines=true,
  columns=fullflexible,
  keepspaces=true,
  upquote=true,
}
\lstdefinestyle{cpshell}{
  language=bash,
  basicstyle=\scriptsize\ttfamily\color{cpInk},
  keywordstyle=\color{cpBlue},
  commentstyle=\color{cpGreen},
  showstringspaces=false,
  breaklines=true,
  columns=fullflexible,
}

% Shared geometry for every box in the document (boxes read as one family).
\tcbset{cpbox/.style={
  enhanced, breakable,
  boxrule=0.5pt, arc=1.5pt,
  left=5pt, right=5pt, top=4pt, bottom=4pt,
  fonttitle=\small\bfseries, coltitle=cpInk,
  attach boxed title to top left={yshift=-2mm,xshift=4mm},
  boxed title style={boxrule=0.3pt, arc=1pt},
}}

% Grey frame + wash fill, shared by all code boxes.
\tcbset{cpcode/.style={cpbox, listing only,
  colback=cpWash, colframe=cpRule,
  boxed title style={colback=cpWash!60!white, colframe=cpRule}}}

\newtcblisting{pycode}[1][Python]{cpcode, listing style=cppython, title={#1}}

\newtcblisting{pycodenum}[1][Python]{cpcode, title={#1},
  listing options={style=cppython, numbers=left,
                   numberstyle=\tiny\color{cpRule}, numbersep=6pt}}

% Untitled variant for two-line snippets where a title box is overkill.
\newtcblisting{pycodeplain}{cpcode, listing style=cppython}

\newtcblisting{shellcode}[1][Shell]{cpcode, listing style=cpshell, title={#1}}

\newtcblisting{jsoncode}[1][JSON]{cpcode, listing style=cpjson, title={#1}}


%% ---------------------------------------------------------------------------
%%  8. LLM prompt boxes
%%
%%  Prompts are data, not prose: they contain braces, underscores, and JSON, so
%%  the default is a VERBATIM box (nothing to escape, copy-pasteable, breaks
%%  across columns and pages):
%%     \begin{sysprompt} ... \end{sysprompt}     system prompt
%%     \begin{userprompt} ... \end{userprompt}   user / augmented prompt
%%     \begin{modelout} ... \end{modelout}       model response
%%     \begin{judgeprompt} ... \end{judgeprompt} LLM-as-judge instructions
%%     \begin{evidence} ... \end{evidence}       retrieved chunk / document
%%
%%  Use the -fmt variants only when a prompt needs typeset emphasis (e.g.
%%  \ph{placeholder} highlighting in the paper's method section); there,
%%  every special character must be escaped by hand.
%% ---------------------------------------------------------------------------

% Placeholder inside a formatted prompt: \ph{break\_date}
\newcommand{\ph}[1]{{\normalfont\itshape\color{cpBlue}\{\{#1\}\}}}

\tcbset{cpprompt/.style={cpbox,
  fontupper=\scriptsize\ttfamily,
  before upper={\parindent0pt\parskip3pt\raggedright\sloppy}}}

% Verbatim prompt text: no language keywords, wrap long lines, mark the wrap.
\lstdefinestyle{cpprompttext}{
  basicstyle=\scriptsize\ttfamily\color{cpInk},
  breaklines=true, columns=fullflexible, keepspaces=true, upquote=true,
  showstringspaces=false,
  postbreak=\mbox{\textcolor{cpRule}{$\hookrightarrow$}\space},
}
\tcbset{cppromptverb/.style={cpprompt, listing only, listing style=cpprompttext}}

\newtcblisting{sysprompt}[1][System prompt]{cppromptverb,
  colback=cpWash, colframe=cpRule,
  boxed title style={colback=cpWash!60!white, colframe=cpRule}, title={#1}}

\newtcblisting{userprompt}[1][User prompt]{cppromptverb,
  colback=cpBlue!4!white, colframe=cpBlue!55!white,
  boxed title style={colback=cpBlue!12!white, colframe=cpBlue!55!white}, title={#1}}

\newtcblisting{modelout}[1][Model output]{cppromptverb,
  colback=cpGreen!4!white, colframe=cpGreen!55!white,
  boxed title style={colback=cpGreen!12!white, colframe=cpGreen!55!white}, title={#1}}

\newtcblisting{judgeprompt}[1][Judge prompt]{cppromptverb,
  colback=cpAmber!5!white, colframe=cpAmber!55!white,
  boxed title style={colback=cpAmber!14!white, colframe=cpAmber!55!white}, title={#1}}

% Formatted (non-verbatim) counterparts — support \ph{} and \textbf{}.
\newtcolorbox{syspromptfmt}[1][System prompt]{cpprompt,
  colback=cpWash, colframe=cpRule,
  boxed title style={colback=cpWash!60!white, colframe=cpRule}, title={#1}}
\newtcolorbox{userpromptfmt}[1][User prompt]{cpprompt,
  colback=cpBlue!4!white, colframe=cpBlue!55!white,
  boxed title style={colback=cpBlue!12!white, colframe=cpBlue!55!white}, title={#1}}
\newtcolorbox{modeloutfmt}[1][Model output]{cpprompt,
  colback=cpGreen!4!white, colframe=cpGreen!55!white,
  boxed title style={colback=cpGreen!12!white, colframe=cpGreen!55!white}, title={#1}}

% Retrieved document / corpus excerpt (RAG sections). Prose, so not verbatim.
\newtcolorbox{evidence}[1][Retrieved document]{cpbox,
  fontupper=\small, before upper={\parindent0pt\parskip3pt},
  colback=cpPlum!4!white, colframe=cpPlum!55!white,
  boxed title style={colback=cpPlum!12!white, colframe=cpPlum!55!white}, title={#1}}

% Qualitative example (input -> output pair discussed in the text).
\newtcolorbox{examplebox}[1][Example]{cpbox,
  fontupper=\small, before upper={\parindent0pt\parskip3pt},
  colback=white, colframe=cpInk!35!white,
  boxed title style={colback=cpInk!8!white, colframe=cpInk!35!white}, title={#1}}


%% ---------------------------------------------------------------------------
%%  9. Domain vocabulary
%%
%%  Every recurring name gets a macro, so a rename is one edit and the paper
%%  can never disagree with itself about a member name or a symbol.
%% ---------------------------------------------------------------------------

% Tool name. REVIEW MODE: the submission is anonymous and the repository link
% must be an anonymized one — check that naming the package does not identify
% the authors (a public PyPI/GitHub project with the same name would).
\newcommand{\pkg}{\textsc{whychange}}

% Ensemble members (exactly the identifiers used in the code and the tables).
\newcommand{\mMosum}{\code{mosum}}
\newcommand{\mPeltL}{\code{pelt\_l2}}
\newcommand{\mPeltN}{\code{pelt\_normal}}
\newcommand{\mKernel}{\code{kernel\_rbf}}
\newcommand{\mBaiPerron}{\code{bai\_perron}}
\newcommand{\mZA}{\code{zivot\_andrews}}
\newcommand{\mProphet}{\code{prophet\_trend}}

% Three-state vote profile: fired / ran-but-silent / not-run.
\newcommand{\vfired}{\textcolor{cpGreen}{\ding{51}}}
\newcommand{\vsilent}{\textcolor{cpRule}{\textendash}}
\newcommand{\vnotrun}{$\emptyset$}
\newcommand{\voteLegend}{%
  \vfired~fired \quad \vsilent~ran, silent \quad \vnotrun~not applicable}

% Notation. Keep these the single source of truth for the maths.
\newcommand{\nobs}{n}                                  % series length
\newcommand{\ts}{\ensuremath{y_{1:\nobs}}}             % the series
\newcommand{\pmem}[1]{\ensuremath{p_{#1}}}             % per-member p-value
\newcommand{\pcomb}{\ensuremath{p_{\mathrm{comb}}}}    % Cauchy-combined p
\newcommand{\epsclust}{\ensuremath{\varepsilon}}       % clustering tolerance
\newcommand{\alphalev}{\ensuremath{\alpha}}            % acceptance level
\newcommand{\blocklen}{\ensuremath{\ell}}              % permutation block length
\newcommand{\nperm}{\ensuremath{B}}                    % number of permutations
\DeclareMathOperator*{\argmax}{arg\,max}
\DeclareMathOperator*{\argmin}{arg\,min}

% Prose shorthands.
\newcommand{\ie}{i.e.,\ }
\newcommand{\eg}{e.g.,\ }
\newcommand{\etal}{et al.\ }


%% ---------------------------------------------------------------------------
%%  10. Drafting aids — MUST be silenced before submission
%%
%%  \draftmodefalse turns every marker below into a no-op, so a stray \todo
%%  cannot reach a reviewer. Flip it in acl_latex.tex.
%% ---------------------------------------------------------------------------

\newif\ifdraftmode
\draftmodetrue            % set \draftmodefalse in acl_latex.tex before submitting

\newcommand{\todo}[1]{\ifdraftmode\textcolor{cpRed}{\textbf{[TODO: #1]}}\fi}
\newcommand{\note}[1]{\ifdraftmode\textcolor{cpBlue}{\textbf{[NOTE: #1]}}\fi}
% A number that still has to be regenerated from the final evaluation run:
% \prov{0.84} is underlined while drafting and plain in the submission build,
% so no placeholder figure can silently survive into the PDF.
\newcommand{\prov}[1]{\ifdraftmode\textcolor{cpRed}{\underline{#1}}\else#1\fi}
% Citation whose bibliographic details are not yet verified against the source
% (EACL screens fabricated references — nothing ships with \unverified live).
\newcommand{\unverified}[1]{\ifdraftmode\textcolor{cpRed}{[verify: #1]}\fi}


% Build-mode test usable in the document body (acl.sty's own flag is internal):
% \ifanonymous ... \else ... \fi  -> true in the [review] build.
\makeatletter
\newif\ifanonymous
\ifacl@anonymize\anonymoustrue\else\anonymousfalse\fi
\makeatother


%% ---------------------------------------------------------------------------
%%  11. Cross-reference names
%% ---------------------------------------------------------------------------

\crefname{algorithm}{Algorithm}{Algorithms}
\Crefname{algorithm}{Algorithm}{Algorithms}
\crefname{section}{Section}{Sections}
\Crefname{section}{Section}{Sections}
\crefname{appsec}{Appendix}{Appendices}
\crefformat{footnote}{#2\footnotemark[#1]#3}


\newcommand{\lfts}{LabelFusion\hbox{-}TS}
\newcommand{\lf}{LabelFusion}
\newcommand{\topone}{top-1}

\title{\lfts{}: Fusing Large Language Models, Transformer Encoders,\\and Financial Time Series for Monetary-Policy Stance Classification}

\author{Internal draft for discussion}

\begin{document}
\maketitle

\begin{abstract}
Financial text is written, read, and interpreted inside a market context, yet financial text classifiers almost universally receive text alone. The surrounding market data, a continuously available and machine-readable signal, is rarely offered to the model as an input modality. The \lf{} architecture combines a prompt-engineered large language model (LLM) with a fine-tuned RoBERTa encoder and was designed so that further input sources can be integrated seamlessly; its authors propose price time series as the natural next modality. We realise this proposal, \lfts{}, and evaluate it on hawkish--dovish classification of Federal Reserve communication, a task where the meaning of a sentence depends on the economic situation it was spoken into. Trained on sentences up to 2015 and tested on 2015--2022, \lfts{} reaches a weighted F1 of 65.0\%, improving over the text-only fusion (62.1\%) and over a zero-shot LLM (64.1\%). The market data contributes exactly where it should: it tells the model which economic regime the words belong to.
\end{abstract}


\section{Introduction}

Most tasks in financial natural language processing are solved by models that read text and nothing else. This is a curious restriction. Every financial document exists inside a market context: analysts who interpret a central-bank statement do so with the rate path, the inflation trajectory, and the equity market in view, and the same sentence can carry opposite meaning in different circumstances. The market context is not an exotic resource. It is public, machine-readable, available at daily resolution for decades, and perfectly aligned in time with any dated document. As an input modality for financial text classifiers it is, nevertheless, largely overlooked: the multimodal literature runs almost entirely in the opposite direction, using text as auxiliary context for price forecasting \citep{tan2025inferring,dolphin2022multimodal}, while classifiers of financial text remain text-only.

\citet{schlee2025labelfusion} introduced \lf{}, which combines the label scores of a prompted LLM with the embeddings of a fine-tuned RoBERTa encoder \citep{liu2019roberta} through a small voting network, and showed that the fusion outperforms either component once enough labelled data is available. The architecture was explicitly designed for the seamless integration of further input sources, and the paper closes by proposing price time series, processed by time-series transformer architectures, as the next modality. This paper carries out that proposal.

We test it where market context should matter most: classifying Federal Reserve communication as hawkish, dovish, or neutral, using the expert-annotated benchmark of \citet{shah2023trillion}. Central-bank language is deliberately hedged. A sentence such as \emph{``the Committee will act as appropriate''} is hawkish during a tightening cycle and dovish during an easing one; the words alone do not decide. Because the interpretation of such language shifts with the economic situation, and because the situation is fully described by public market series, this task is a natural probe for the modality. We evaluate the way the system would be used: trained on the past, tested on the future.

Our contributions are:
\begin{enumerate}
\item \lfts{}, which extends \lf{} with a third expert: a small time-series transformer that reads the market context of each document, realising the extension proposed in the original work.
\item A time-honest evaluation of monetary-policy stance classification: models are trained on sentences up to 2015 and tested on 2015--2022, so no knowledge of the test years can leak into training.
\item Evidence for the overarching claim: adding the market modality improves the fusion from 62.1\% to 65.0\% weighted F1 and lifts it past a modern zero-shot LLM, in a regime with only ${\sim}1{,}300$ labelled training sentences.
\end{enumerate}


\section{Related Work}

\paragraph{\lf{}.}
Given a text $x$, \lf{} \citep{schlee2025labelfusion} obtains label scores $z = f_{\mathrm{LLM}}(x)$ from a prompted LLM and a contextual embedding $h = f_{\mathrm{RB}}(x)$ from a fine-tuned RoBERTa encoder, and fuses them with a small MLP, $\hat{y} = f_{\mathrm{MLP}}([\,h;z\,])$. On Reuters-21578 the fusion beats both components at full training data, while the prompted LLM dominates in the low-data regime. We keep this design unchanged and add one expert.

\paragraph{Monetary-policy stance classification.}
\citet{shah2023trillion} built the benchmark we use: sentences from FOMC meeting minutes, press-conference transcripts, and speeches (1996--2022), each labelled hawkish, dovish, or neutral by trained annotators, together with baselines from rule-based systems to fine-tuned RoBERTa-large and zero-shot ChatGPT. Their own market analysis runs downstream of classification, relating predicted stance to asset prices. To our knowledge, no published work uses market data as an \emph{input} to stance classification; the connection between central-bank communication and markets is otherwise well documented \citep{tetlock2007giving,barber2008glitters}.

\paragraph{Time series and text.}
Purpose-built series encoders such as PatchTST \citep{nie2023patchtst} treat a time series as a sequence of patches, which suits short fixed windows. The multimodal time-series literature predominantly uses text to help forecast series; the reverse direction, series helping to classify text, is the gap this paper addresses.


\section{Data and Task}
\label{sec:data}

\paragraph{Benchmark.}
The dataset of \citet{shah2023trillion} contains 2{,}479 expert-annotated sentences (hawkish / dovish / neutral) drawn from three types of Federal Reserve communication between 1996 and 2022. The task is single-label, three-class classification of one sentence at a time; the benchmark's headline metric is weighted F1.

\paragraph{Dating the sentences.}
Connecting a sentence to market data requires knowing \emph{when} it was said, and the released sentences carry only a year. We therefore matched every sentence back to its dated source document (all source documents are public and included in the authors' repository) by exact and containment text matching, recovering an unambiguous document date for 1{,}692 sentences (73\%); the remainder, mostly boilerplate that recurs across documents, is excluded. Details in Appendix~\ref{app:dating}.

\paragraph{Market context.}
For each sentence we build a window of six daily series over the 126 business days ending the day \emph{before} its document date: the federal funds rate, the 2-year Treasury yield, the yield-curve slope (10y minus 2y), the log equity index, CPI year-over-year inflation, and the unemployment rate, each expressed as its change since the window start. All series are public (FRED). Two safeguards keep the window honest: it ends before the document date, and the macro releases enter with a 45-day lag, so the window contains only what markets actually knew at the time.

\paragraph{Evaluation protocol.}
We split by time: training on the 1{,}274 sentences up to September 2015 (with the most recent 15\% held out for model selection) and testing on the 418 sentences from 2015--2022. This mirrors deployment, where a model trained on past communication must read future statements, and it prevents a subtle leak that a random split would allow: sentences from the same era appearing on both sides and revealing the regime.


\section{\lfts{}}
\label{sec:model}

\begin{figure}[t]
  \centering
  \figph{\columnwidth}{4.2cm}{architecture diagram}
  \caption{\lfts{}. Three experts are trained independently: a fine-tuned RoBERTa encoder reads the sentence; a prompted LLM votes on the sentence; a small time-series transformer reads the market window of the sentence's document. A voting MLP, trained separately with the experts frozen, combines the RoBERTa embedding, the LLM vote, and the market embedding into the final stance prediction, exactly as in \lf{} with one additional input.}
  \label{fig:arch}
\end{figure}

\lfts{} keeps the recipe of \lf{}: each expert is trained on the task individually, then frozen, and a small voting network is trained separately on their outputs (Figure~\ref{fig:arch}).

\paragraph{Text expert.}
RoBERTa-large \citep{liu2019roberta} is fine-tuned on the training sentences; its CLS embedding $h \in \mathbb{R}^{1024}$ represents the sentence, as in \lf{}.

\paragraph{LLM expert.}
An instruction-tuned open LLM\footnote{\texttt{gemma4:31b} at temperature 0. We use verbatim the classification prompt that \citet{shah2023trillion} used for ChatGPT, so this expert is exactly the benchmark's LLM baseline with a current model.} classifies each sentence zero-shot; its vote is encoded as $z \in \{0,1\}^{3}$.

\paragraph{Market expert.}
A small patch transformer in the spirit of PatchTST \citep{nie2023patchtst} reads the market window: the $126 \times 6$ window is cut into 18 week-long patches, each linearly embedded; four transformer layers (width 64, ${\sim}0.4$M parameters) exchange information across the weeks, and the averaged output is the market embedding $m \in \mathbb{R}^{64}$. The expert is trained in two stages, both independent of the other experts: first on six decades of historical windows to predict whether the policy rate changes within the next month, which teaches it what rate cycles look like before it ever sees a labelled sentence, and then on the training sentences (window $\rightarrow$ stance label).

\paragraph{Fusion.}
With all experts frozen, a voting MLP with one hidden layer combines the three views of a sentence,
\begin{equation}
\hat{y} \;=\; f_{\mathrm{MLP}}\big(\big[\,h;\; z;\; m\,\big]\big) \;\in\; [0,1]^{3},
\label{eq:fusion}
\end{equation}
which is \lf{}'s fusion with one added input. Removing $m$ recovers \lf{} exactly; this is the comparison that isolates what the market modality contributes.


\section{Results}
\label{sec:results}

\begin{table}[t]
  \caption{Stance classification on the time-split test set (418 sentences from 2015--2022; models trained on sentences up to 2015). Weighted F1 and accuracy in percent, following the benchmark's metric.}
  \label{tab:main}
  \centering
  \begin{fitbox}
  \tabfont
  \begin{tabular}{@{}lcc@{}}
    \toprule
    & wF1 (\%) & Acc.\ (\%) \\
    \midrule
    majority class                     & 26.9 & 44.0 \\
    market expert alone                & 38.7 & 39.2 \\
    \midrule
    \lf{} (text + LLM)                 & 62.1 & 62.0 \\
    RoBERTa-large, fine-tuned          & 62.6 & 62.4 \\
    LLM zero-shot                      & 64.1 & 64.6 \\
    \lfts{} (ours)                     & \best{65.0} & \best{64.8} \\
    \bottomrule
  \end{tabular}
  \end{fitbox}
\end{table}

Table~\ref{tab:main} reports the time-split evaluation. The market expert alone reaches 38.7\% weighted F1: market data by itself carries a real but weak stance signal, well above the majority baseline and far below any text model. The fine-tuned encoder reaches 62.6\% and the text-only fusion 62.1\%; the zero-shot LLM, with no training data at all, reaches 64.1\% and outperforms both. Adding the market expert changes the picture: \lfts{} reaches 65.0\%, a gain of $+2.9$ points over the text-only fusion and the only configuration that surpasses the zero-shot LLM.

Two aspects of this result are worth stating precisely. First, the entire improvement comes from adding one modality: architecture, experts, and training recipe are otherwise identical between \lf{} and \lfts{}. Second, the pattern reproduces the central observation of \citet{schlee2025labelfusion} under harder conditions. In their experiments, the fusion needed roughly 80\% of a much larger training set before it overtook the prompted LLM. Here, with only ${\sim}1{,}300$ labelled sentences and a test set drawn from years the models never saw, the text-only fusion still trails the LLM, exactly as their low-data results predict, and it is the market modality that closes the gap.


\section{Discussion}
\label{sec:discussion}

\paragraph{Why market context helps here.}
The mechanism is visible in the data before any model is trained: in documents preceded by six months of rising rates, 65\% of the stance-bearing sentences in the training years are hawkish; after falling rates, 25\%. The economic situation shifts the meaning of Fed language, and hedged sentences inherit their stance from the regime they are spoken into. The market window hands the classifier exactly this missing variable. Consistent with that reading, the market expert is nearly useless alone (38.7\%) but valuable in combination: its embedding lets the fusion reinterpret the text, which is why the contribution appears only when text and market enter the model together.

\paragraph{Relation to the zero-shot LLM.}
The comparison with the zero-shot LLM understates our case in one respect: the LLM was trained on data that includes the 2015--2022 test years and public commentary about them, whereas the fine-tuned models saw nothing after 2015. That the fused system nevertheless edges past it suggests the market modality contributes information that even era-exposed pretraining does not supply on demand.

\paragraph{Limitations.}
We report single training runs, as does the benchmark and \citet{schlee2025labelfusion}; a seed-variance analysis is in preparation. The study covers one task, one benchmark, and a test span of seven years dominated by two unusual regimes (the zero-rate era and the 2022 tightening). Dating recovered 73\% of the benchmark's sentences, so our split is not identical to the official one. And the market expert's standalone weakness means the modality is a complement to text, not a substitute; nothing here suggests market data alone classifies financial text.


\section{Conclusion}

Financial text classifiers routinely ignore an input that every human analyst consults: the market context. We added it to an existing fusion architecture as one further expert, exactly as that architecture's authors proposed, and evaluated the result on central-bank stance classification under a train-on-the-past, test-on-the-future protocol. The market modality improved the fusion by $2.9$ points of weighted F1 and lifted it past a zero-shot LLM in the low-label regime where text-only fusion falls behind. The broader point is simple: for financial tasks, the time series surrounding a document is cheap, public, perfectly time-aligned, and informative, and there is little reason to keep leaving it out.

% Internal draft: further results (seed analysis, ablations over ~100 TS-expert and fusion
% configurations, additional evaluation protocols) exist and will be added before submission.

\bibliography{refs}

\appendix

\section{Dating Procedure}
\label{app:dating}

The benchmark releases sentences with labels and a year, but not the document date needed to attach market data. All source documents are public and included in the benchmark's repository as date-named files. We match each sentence to its document by exact text match against the per-document sentence files and, where that fails, by containment matching of the normalised sentence against the raw document text. Sentences matching documents on more than one date (recurring boilerplate) and sentences matching no document are excluded. This yields 1{,}692 uniquely dated sentences of 2{,}312 unique sentences (73\%), with label proportions close to the full benchmark (dovish 28\%, hawkish 25\%, neutral 47\%).

\section{Market Window Details}
\label{app:window}

Series (all FRED): federal funds effective rate (DFF), 2-year Treasury yield (DGS2), yield-curve slope (DGS10 $-$ DGS2), Nasdaq Composite as log index (NASDAQCOM), CPI year-over-year from CPIAUCSL, and the unemployment rate (UNRATE). Windows span 126 business days and end the day before the document date. CPI and unemployment enter with a 45-day publication lag. Each channel is expressed as its change since the window start and standardised with statistics computed on the training split only.

\section{LLM Prompt}
\label{app:prompt}

The LLM expert uses, verbatim, the prompt published by \citet{shah2023trillion} for their ChatGPT baseline: \emph{``Discard all the previous instructions. Behave like you are an expert sentence classifier. Classify the following sentence from FOMC into `HAWKISH', `DOVISH', or `NEUTRAL' class. Label `HAWKISH' if it is corresponding to tightening of the monetary policy, `DOVISH' if it is corresponding to easing of the monetary policy, or `NEUTRAL' if the stance is neutral. Provide the label in the first line and provide a short explanation in the second line. The sentence: [sentence]''} The first line of the reply is parsed as the vote.

\section{Training Details}
\label{app:training}

Text expert: \texttt{roberta-large}, maximum length 128 tokens, AdamW with learning rate $10^{-5}$, class-weighted cross-entropy, batch size 16 with gradient accumulation 2, up to 6 epochs with epoch selection on the held-out validation tail, mixed precision, single T4 GPU (${\sim}8$ minutes). Market expert: patch length 7 (18 patches), width 64, 4 layers, 8 heads, dropout 0.2; pretraining on windows sampled every second business day 1962--2015 with a 3-class rate-change target (30-day horizon, $\pm 10$bp threshold); task training with Adam at $3 \times 10^{-4}$, up to 40 epochs with validation selection. Voting MLP: one hidden layer of 64, trained on the concatenated expert outputs with the experts frozen. The market expert and fusion train in minutes on CPU.

\end{document}
